\section{Diskussion}
%Einleitung in das Kapitel
Um einen umfassenden Überblick über die erzielten Ergebnisse der vorliegenden Arbeit zu geben, werden im Folgenden die zentralen Erkenntnisse zusammengefasst und kritisch reflektiert. Dabei wird insbesondere auf die identifizierten Effizienzpotenziale des entwickelten \ac{KI}-Agenten im Vergleich zum herkömmlichen \ac{ERP}-Standardprozess eingegangen und deren Bedeutung im Kontext der bestehenden Forschungsliteratur diskutiert. Zudem werden die Limitationen der Arbeit dargestellt, um die Ergebnisse passend einordnen zu können.
%Unterkapitel
\subsection{Zusammenfassung der Ergebnisse}

Im Vergleich zum \ac{ERP}-Standardprozess ermöglicht der entwickelte \ac{KI}-Agent eine Reduktion der vom Anwender aktiv auszuführenden Prozessschritte. Dies gilt sowohl für die Anlage des Angebotskopfes als auch für die Anlage der Angebotspositionen. Insbesondere bei der Artikelsuche und der Preisfindung werden manuelle Navigationsschritte sowie wiederholte wechsel zwischen verschiedenen Systemmasken vollständig vermieden. Darüber hinaus ist es durch die dialogbasierte Interaktion möglich, mehrere Eingaben gebündelt zu erfassen und so die Anzahl der erforderlichen Systeminteraktionen weiter zu reduzieren. 

Darüber hinaus zeigt sich ein Effizienzpotenzial in der inhaltlichen Unterstützung der Anwender. Während das \ac{ERP}-System im Standardprozess primär systemseitige Pflichtfelder validiert, überprüft der \ac{KI}-Agent zusätzlich unternehmensspezifisch relevante Eingaben. Dadurch können unvollständige oder fehlerhafte Angebotsdaten bereits während der Angebotsanlage erkannt werden. Dieses Potenzial wirkt sich nicht nur auf die unmittelbare Prozessdurchführung aus, sondern kann auch zu einer Reduktion von Nacharbeit im weiteren Vertriebsprozess beitragen.

%Die Ergebnisse stehen im Einklang mit dem im Stand der Forschung beschriebenen Einsatz von \ac{KI}-basierten Assistenzsystemen in \ac{ERP}-Systemen. Dort wird insbesondere hervorgehoben, dass der Mehrwert von \ac{KI} weniger in der vollständigen Automatisierung komplexer Geschäftsprozesse liegt, sondern vielmehr in der situationsabhängigen Unterstützung von Anwendern bei wissensintensiven und variablen Tätigkeiten. Der entwickelte \ac{KI}-Agent bestätigt diese Erkenntnisse, indem er keine autonomen Entscheidungen trifft, sondern den Anwender durch strukturierte Dialoge, Vorschläge und Validierungen unterstützt.

Insgesamt verdeutlichen die Ergebnisse, dass der Einsatz eines \ac{KI}-Agenten im Kontext der Angebotserstellung vor allem dort sinnvoll ist, wo Prozesse durch hohe Variantenvielfalt, unvollständige Informationen oder manuelle Such- und Interpretationsaufgaben geprägt sind.

\subsection{Limitationen der Arbeit}

Allerdings sind bei der Interpretation der Ergebnisse auch einige Limitationen zu berücksichtigen. Vorab ist dabei klarzustellen, dass es sich bei den in dieser Arbeit identifizierten Ergebnissen bewusst um Effizienzpotenziale und nicht um empirisch nachgewiesene Effizienzsteigerungen handelt. Ziel der Arbeit war es, potenzielle Verbesserungen auf Prozessebene zu identifizieren und zu analysieren, nicht jedoch eine quantitative Bewertung tatsächlicher Effizienzgewinne vorzunehmen. Die identifizierten Effizienzpotenziale beruhen also auf einer strukturellen Analyse der Prozessabläufe und der Reduktion manueller Prozessschritte, nicht jedoch auf einer quantitativen Messung der tatsächlichen Zeitersparnis.

Darüber hinaus ist zu berücksichtigen, dass, wie bereits in der Literaturrecherche beschrieben, die Effektivität von \ac{KI}-Agenten maßgeblich von der Qualität und Verfügbarkeit der zugrundeliegenden Daten abhängt \cite[\pagef 1035]{gujar_role_2025}. In der vorliegenden Arbeit wurde eine demonstrative Datenbasis verwendet, die über eine konsistente und vollständige Datenbasis verfügt, jedoch die Komplexität und potenzielle Inkonsistenzen realer Unternehmensdaten nur eingeschränkt abbildet. Also wurde davon ausgegangen, dass alle relevanten Stammdaten, wie Kunden- und Artikelinformationen, sowie historische Angebotsdaten vollständig und korrekt vorliegen, weshalb der \ac{KI}-Agent nicht mit Problemen wie fehlenden Einträgen, Dateninkonsistenzen oder unklaren Datenformaten umgehen musste.

Neben der Datenqualität wurden weitere Aspekte eines realen Produktivbetriebs nicht untersucht. Hierzu zählen insbesondere Performanceeffekte unter hoher Systemlast sowie Aspekte wie langfristige Stabilität, Wartbarkeit und Skalierbarkeit, aber auch organisatorische Faktoren wie Nutzerakzeptanz, Vertrauen in \ac{KI}-gestützte Unterstützung oder Veränderungen etablierter Arbeitsweisen. Diese Faktoren können einen erheblichen Einfluss auf die tatsächliche Realisierung der identifizierten Effizienzpotenziale des \ac{KI}-Agenten haben.

Zudem ist zu beachten, dass sich die Arbeit auf einen spezifischen Anwendungsfall innerhalb des Vertriebsprozesses von \ac{D365FnSCM} konzentriert. Die Übertragbarkeit der Ergebnisse auf andere Geschäftsprozesse oder Branchen ist daher nur eingeschränkt gegeben. Zwar bleibt das grundlegende Ergebnis bestehen, dass sich durch den Einsatz von \ac{KI}-Agenten Effizienzpotenziale realisieren lassen können, deren Ausprägung kann jedoch
je nach Prozesscharakteristik, Organisationsstruktur und Systemlandschaft deutlich variieren.

Schließlich erfolgte die Entwicklung und Integration des \ac{KI}-Agenten vollständig innerhalb des Microsoft-Ökosystems. Die dadurch vorhandene enge Verbindung von \ac{D365FnSCM}, Dataverse und Microsoft Copilot Studio erleichtert die technische Umsetzung sowie den Zugriff auf \ac{ERP}-Daten erheblich. Die identifizierten Effizienzpotenziale sind daher auch vor dem Hintergrund dieser integrierten Plattform zu betrachten. In anderen \ac{ERP}-Systemen, bei denen vergleichbare Integrationsmechanismen nicht oder nur eingeschränkt verfügbar sind, können sich sowohl der Entwicklungsaufwand als auch die realisierbaren Effizienzpotenziale unterschiedlich darstellen.